\subsection[Usage of the BMA package]{Usage of the \pkg{BMA} package}\label{sec:usage_of_the_BMA package}
\subsubsection{The GUI way}
Once you start \pkg{gretl}, you must open a data file and then you can
load the relevant \pkg{BMA} package from the \pkg{gretl} server. In the main
window, go to Tools $>$ Function packages $>$ On server heading. By
selecting \pkg{BMA}, you will open a window similar to the one shown in
Figure~\ref{fig:Main window for BMA}:

\begin{figure}[ht!]
	\centering
	\includegraphics[width=0.7\columnwidth, keepaspectratio]{bma_main_window}
	\caption{Main window for \pkg{BMA}.}
	\label{fig:Main window for BMA}
\end{figure}

According to Figure~\ref{fig:Main window for BMA}, we can specify the following
entries in the GUI \pkg{BMA} window\label{text:module_parameters}
\begin{itemize}
	\item \textbf{Dependent variable} -- The dependent variable.
	\item \textbf{List of independent variables}
		-- Loading variables from the database, which must have been
		opened previously.  Note that by
		default we assume that you want to estimate an intercept; therefore,
		a constant is implicitly included in the list of the variables.
	\item \textbf{Model prior} -- Indicates the choice of model prior. One
		can employ the binomial model prior or the binomial-beta model
		prior. Note that the uniform model prior is a special case of the
		binomial model prior. Therefore, in fact, our package allows for
		three types of priors.
	\item \textbf{Prior average model size} -- Specifies the prior
		expected model size $\E(\Xi)$.  The expected model size may range
		from 0 to \textit{K}. The default value equals 1.  It is the
		smallest reasonable prior expected model size. It means that we
		penalize large models and we assign high probability to small ones
		following Occam's razor philosophy.
	
		Note that for the binomial model prior and with $\E(\Xi)=0.5K$ one
		can define the uniform prior on the model space. For example, for
		$K=10$ regressors we can define the uniform prior with
		$\E(\Xi)=5$. For $\E(\Xi) < 5$ we assign high probability to small
		models.  The smaller the prior expected model size $\E(\Xi)$, the
		less probable are larger models.
	\item \textbf{Significance level for the initial model} -- Defines the
		significance level which is used to build the initial model. An
		explanatory variable enters the initial model if its $p$~value is
		less than the significance level. If the significance level equals
		1, the initial model will be randomly chosen (with equal
		probability) from all available models. Note that if all available
		explanatory variables enter the initial model, you will get the
		following \pkg{gretl} error message ``No independent variables were
		omitted''.
	\item \textbf{Number of top ranked models} -- Specifies the number of
		best models for which detailed information is stored.
	\item \textbf{g-prior type} -- One can choose between four Zellner's
		g-priors for the regression coefficients. Choices include: benchmark
		prior, unit information prior, risk inflation criterion, Hannan and
		Quinn prior, root of g-UIP.
	\item \textbf{Jointness analysis} -- If \code{"None"} (the default),
		the jointness analysis is omitted. Alternatively, one can choose
		one of the jointness measures discussed in section~\ref{sec:jointness_measures}.
	\item \textbf{Number of out-of-sample forecasts} -- Defines the total
		number of out-of-sample forecasts of the dependent variable.
	\item \textbf{Total number of replications} -- Defines the total
		number of iteration draws to be sampled.
	\item \textbf{Percentage of burn-in draws} -- Specifies the number of
		burn-in replications, calculated as the percentage of the total
		number of iteration draws.
	\item \textbf{Verbosity} -- An integer value of 1 or 2; the default is
		1, which allows to see the basic Bayesian model averaging
		results. If Verbosity equals 2, a more detailed description of the
		analysis is provided (initial model, speed of convergence,
		estimation results for top ranked models).
	\item \textbf{Show overall progress} -- The output is flushed every 10\%
		of MC$^3$ and the progress is printed. Note: this consumes a lot of CPU time.
\end{itemize}     

\subsubsection{Script}
The \pkg{BMA} package can also be used inside \proglang{Hansl} scripts.
A very minimal code could be as follows:
%
\begin{lstlisting}
 open greene9_1.gdt
 include BMA.gfn
  
 list greene = capital labor nfirm
 
 BMA(valadd, greene)
\end{lstlisting}
The above code triggers BMA analysis with default option
described on page~\pageref{sec:default_option_bundle}.
On contrary, the corresponding ''big call'', i.e. with
setting all possible parameters, would be as follows:
\begin{lstlisting}
bundle opts = defbundle("acc_type", 1, "av_model_size", 1.5, \
 "alpha", 0.6, "l_rank", 4, "g_type", 1, "do_joint", 0, "h_predict", 0, \
 "Nrep", 1000, "burn", 10, "verbosity", 2, "show_progress", 1)
 
 BMA(valadd, greene, opts)
\end{lstlisting}

or alternatively with the new syntax for creating a bundle
on-the-fly:
\begin{lstlisting}
 BMA(valadd, greene, _(acc_type=1, av_model_size=1.5, alpha=0.6, \
  l_rank=4, g_type=1, do_joint=0, h_predict=0, Nrep=1000, \
  burn=10, verbosity=2, show_progress=1))
\end{lstlisting}
The above example consists of three blocks.

The first block defines just the opening of the so-called \code{greene9_1}
dataset, which is bundled in every standard \pkg{gretl}
installation. This dataset contains the cross-sectional data on the
manufacturing of transportation equipment presented in
\cite{greene_Econometric_A}%
\footnote{This dataset is also available at link
	\url{http://people.stern.nyu.edu/wgreene/Text/Edition7/TableF7-2.txt}.}
and it consists of the following variables: valadd -- output, capital
-- capital input, labor -- labor input and nfirm -- number of firms.

The second block is the definition of the \code{greene} list which
contains indenpendent variables available in the \code{greene9_1} dataset.

The third block contains the specification of the Bayesian model
averaging analysis: binomial model prior, prior average model size set
to $1.5$ (which means that we have the uniform model prior), significance
level for the initial model set to $0.6$, 4 top ranked models,
Benchmark g-prior, without jointness analysis, without out-of-sample
forecasts, 1000 replications with 10\% burn-in draws and overall
progress is printed.

Suppose we want to run the same analysis but following Occam's razor
philosophy.  The simplest way to do this, is to set the prior average
model size to 1.  The code would be as follows:
%
\begin{lstlisting}
 BMA(valadd, greene, _(av_model_size=1))
\end{lstlisting}
%
In this case, we penalize large models by assigning high probabilities
to small ones. Apart from the binomial model prior, we can use the
binomial-beta prior. To do this, we have to set the model prior to 2.
The code would be as follows:
%
\begin{lstlisting}
 BMA(valadd, greene, _(av_model_size=1, acc_type=2))
\end{lstlisting}
%
Suppose we want to set the g-prior to risk inflation criterion,
select the uniform model prior, do jointness analysis with the
modified Yule's~Q (YQM) measure, compute forecasts for 3
observations, print just basic output (verbosity set to 1) and
switch off progress printing. The code would then be as follows:
%
\begin{lstlisting}
 BMA(valadd, greene, _(g_type=3, do_joint=7, h_predict=3, verbosity=1, show_progress=0))
\end{lstlisting}
%
Finally, if we want to save results of the above Bayesian model
averaging analysis in the the bundle \code{results_b}, the code
should be as follows
%
\begin{lstlisting}
 results_b = BMA(valadd, greene, _(g_type=3, do_joint=7, \
  h_predict=3, verbosity=1, show_progress=0))
\end{lstlisting}
%
Suppose we want to compute Doppelhofer-Weeks (DW2) Measure instead of
already computed Yule's~Q (YQM). Since the type of jointness measure
is controlled by 'do\_joint' field in the result bundle, the only thing
we have to do is to set the value of 'do\_joint' to 2 and print BMA
results again. The code should be as follows:

\begin{lstlisting}
 results_b.do_joint = 2
 BMA_Print(&results_b)
\end{lstlisting}
%
If we need backwards compatibility with BMA\_package prior to version 2
we can use the BMA\_OldStyle function to run the old style script.

\begin{lstlisting}
 list greene_old = valadd capital labor nfirm
 BMA_OldStyle(greene_old, 1, 1.5, 0.6, 4, 1, 1, 3, 1000, 10, 2)
\end{lstlisting}

\subsubsection{Outputs}\label{sec:outputs}
If you select the appropriate entries in the GUI \pkg{BMA} window, our
package returns: PIP, Mean, Std.Dev., Cond.Mean, Cond.Std.Dev.
Furthermore, the package also returns the predictive results of the
dependent variable (Mean and Std.Dev.).

Let us consider the data used in \cite{FernandezLeySteel2001b}
(denoted FLS hereafter). This data comprises the information about 72
countries and 41 potential growth determinants for the period 1960 to
1992%
\footnote{The dataset is publicly available on the \emph{Journal of
		Applied Econometrics} online data archive
	(\url{http://qed.econ.queensu.ca/jae/2001-v16.5/fernandez-ley-steel/}).
	The reported chain (3000000 replications with 33\% burn-in draws)
	took about 37.5 minutes of CPU time on a PC with AMD Phenom II X6
	1100T CPU, 6.0~Gb of RAM running under Debian GNU/Linux. We used
	\pkg{gretl}~2022b-git compiled by GCC~11.2.0. The seed for RNG was set
	to 1000000. The dependent variable is the growth rate from
	1960--1992.}. The initial estimation sample consists of indices
$i=1,\ldots,69$ and the out-of-sample evaluation indices are from 70
to 72.

For example, for the FLS data, the following estimates should appear
%
\begin{CodeChunk}
	\begin{CodeOutput}
Posterior moments (unconditional and conditional on inclusion):
PIP         Mean     Std.Dev.    Cond.Mean   Cond.Std.Dev        
GDPsh560    0.998378    -1.444579     0.305391    -1.446927       0.300031
Confuncious 0.890878     0.427081     0.189631     0.479393       0.123638
Life_Exp    0.932297     0.988063     0.404005     1.059816       0.314688
Equip_Inv   0.904971     0.549653     0.250217     0.607370       0.184734
SubSahara   0.654270    -0.431560     0.369153    -0.659606       0.240550
Muslim      0.363769     0.130498     0.197138     0.358739       0.157977
Rule_of_Law 0.416318     0.206002     0.271416     0.494818       0.184491
Yrs_Open    0.547074     0.312026     0.318952     0.570354       0.196508
...         ...          ...          ...          ...            ...
----------------------------------
Posterior probability of models:
Model 1:    0.007899
Model 2:    0.005294
Model 3:    0.004204
Model 4:    0.003960
Model 5:    0.003378
...         ...
Total probability of the models in ranking (numerical): 0.058997
Correlation coefficient between the analytical
and numerical probabilities of the above models: 0.997518
----------------------------------
Predictive results:
Growth       Mean   Std.Dev.
Obs.70    2.714600   1.590657   1.279354
Obs.71    2.067600   2.006087   0.941517
Obs.72    1.889700   2.465582   0.981258
	\end{CodeOutput}
\end{CodeChunk}
%
The \pkg{BMA} estimate function accepts a scalar, which sets the verbosity
of the output. Its default value is 1, which causes the estimation
output to be printed out. A value of 2 forces the \pkg{BMA} function to
print out all the details of the estimation. You can print out the
above-mentioned results and additionally the following information:
the total CPU time, type of model prior, prior average model size,
significance level for the initial model, type of g-prior, total
number of iterations and, finally, the number of burn-in
draws. Moreover, the \pkg{BMA} estimate function produces the estimation
results for the initial and top ranked models. The jointness analysis
is inactive by default. If it is active, you will get: posterior joint
probability of explanatory variables, jointness statistic as given in
Equations~\ref{DW_jointness1}--\ref{YqM_jointness} and classification
of jointness measures.  The jointness analysis for the previous
example should look like this:
%
\begin{CodeChunk}
	\begin{CodeOutput}
Posterior joint probability of variables:
GDPsh560    Confuncious   Life_Exp  Equip_Inv ...  
GDPsh560    0.000000    0.889256   0.932182   0.903349 ...
Confuncious 0.000000    0.000000   0.826208   0.816956 ...
Life_Exp    0.000000    0.000000   0.000000   0.839068 ...
Equip_Inv   0.000000    0.000000   0.000000   0.000000 ...
...         ...         ...        ...        ... ...
Jointness statistics (Ley-Steel Measure (LS2)):
GDPsh560    Confuncious   Life_Exp  Equip_Inv ...
GDPsh560    0.000000    2.083161   2.643181   2.235000 ...
Confuncious 0.000000    0.000000   1.576599   1.618373 ...
Life_Exp    0.000000    0.000000   0.000000   1.662559 ...
Equip_Inv   0.000000    0.000000   0.000000   0.000000 ...  
...         ...         ...        ...        ... ... 
Strong substitutes:
Rev_Coup,Area         -4.516140
Rev_Coup,Publ_Edu_pct -4.459739
...             ...
Significant substitutes:   
GDPsh560,Hindu        -1.997989
Mining,War_Dummy      -1.989541
...                   ...   
Significant complements:
Confuncious,Life_Exp   1.576599
Confuncious,Equip_Inv  1.618373
...                    ...
Strong complements:   
DPsh560,Confuncious    2.083161
GDPsh560,Equip_Inv     2.235000
...                    ...  
	\end{CodeOutput}
\end{CodeChunk} 
